\subsection{Set Up for Resource Theory}
我们希望有一个理论框架能够描述量子系统的一些「特质」。如同 coherence以及entanglement等等。Resource
Theory正是给出了这样的一个框架。而被描述的这些「特质」我们称之为资源（resource）。

\bigskip
\hlr{Resource Theory 的基本要素}

对于一个能够描述资源的Resource Theory，我们需要给出以下几个要素：
\begin{itemize}
  \item \textbf{Free
    Operations}：一系列操作，保证这些操作并不会产生新的资源。并且注意\textbf{Trace
    Operation}必然是Free Operation的一部分，因为丢弃一个子系统不可能产生资源。
  \item \textbf{Free States and Resource States}：free states是只能通过free
    operations产生的状态，表征没有资源的状态；resource states是不能通过free operations产生的状态，表征有资源的状态。
\end{itemize}
我们通过严格定义「什么态没有资源」以及「什么操作不会产生资源」来建立一个Resource Theory。

\bigskip
\hlr{Simulation of Resource}

我们可以通过一个resource state和 free operation来模拟一个产生资源的操作。如果存在一个free
operation $ \mathcal{F} $和一个resource state $ \sigma $使得对于任意的state $ \rho $都有：
\begin{align}
  \mathcal{E}(\rho) = \operatorname{Tr}_{A}\left[\mathcal{F}(\rho
  \otimes \sigma)\right],
\end{align}
那么我们称$ \mathcal{E} $可以被$ \sigma $和free operation $ \mathcal{F} $模拟。
这也告诉我们，我们可以通过resource state来实现一个non-free的operation。

\bigskip
\hlr{Interconversion}

对于两个state $ \rho $和$ \sigma $，如果存在一个free operation $ \mathcal{E} $使得$
\mathcal{E}(\rho) = \sigma $，我们称之为$ \rho \to \sigma $或者$ \rho \prec
\sigma $是一个interconversion。这个操作必然是可以传递的：
\begin{align}
  \rho_{1}\prec\rho_{2}\mathrm{~and~}\rho_{2}\prec\rho_{3}\Longrightarrow\rho_{1}\prec\rho_{3}.
\end{align}
这个要求其实就是在说，free operation给出了state 之间的一个partial ordering。数学上定义为:
\defi{
  Partial Ordering

  对于一个集合$ S $，如果存在一个关系$ \preceq $满足：
  \begin{itemize}
    \item Reflexive: 对于任意$ a \in S $都有$ a \preceq a $
    \item Antisymmetric: 对于任意$ a,b \in S $如果$ a \preceq b $并且$ b
      \preceq a $那么$ a = b $
    \item Transitive: 对于任意$ a,b,c \in S $如果$ a \preceq b $并且$ b
      \preceq c $那么$ a \preceq c $
  \end{itemize}
}
对于第一个要求显然有Identity Operator是一个free operation使得$ \rho \to \rho
$成立。第二个要求在合理定义等价类的情况下也成立。第三个要求就是上面的传递性。

我们可以证明所有free state都是可以互相转换的。因为Trace是free operation，我们永远可以通过free
operation trace掉一个state然后再通过free operation准备另一个free state来实现两个free state之间的转换。

\bigskip
\hlr{Monotone}

在Resource Theory中，我们希望能够定义一个函数来度量资源的多少，这个函数被称为monotone。对于一个monotone $
M(\rho) $，我们要求其满足，Free Operations不会增加monotone的值：
\begin{align}
  M(\rho)\geq M(F(\rho)).
\end{align}

\subsection{Resource Theory of Entanglement}

\subsubsection{Free Operations and Free States}

\bigskip
\hlr{Free Operations for Entanglement}

下面我们希望建立一个描述entanglement的Resource Theory。我们需要知道Free
Operation是什么我们发现有三种Free Operations：
\begin{itemize}
  \item \textbf{Local Operation} 也就是说$ \mathcal{E}( * ) =
    \mathcal{E}_A \otimes \mathcal{E}_B ( * ) $，其中$ \mathcal{E}_A,
    \mathcal{E}_B $都是作用在各自子系统的Unitary操作。
  \item \textbf{Trace:} 也就是说$ \mathcal{E}( * ) = \operatorname{Tr}_A(
    * ) $或者$ \mathcal{E}( * ) = \operatorname{Tr}_B( * ) $，也就是丢弃一个子系统。
  \item \textbf{Classical Communication:}
    也就是对于系统的经典信息的Relabling，我们可以通过下面的定义严格描述：
    \begin{align}
      C(\rho) = \sum_n \ketbra{n}{n}_B \bra{n} \rho \ket{n}_A \quad
      \rho \in \mathcal{H}_A \quad C(\rho) \in \mathcal{H}_B
    \end{align}
    就相当于把A所包含的经典信息传递给B。
\end{itemize}
一个很经典的例子就是，我们可以通过A系统得到的结果来对于B系统进行条件操作，这是一个经典信息的传输，我们完全可以通过把A的结果写在纸上然后递给操作B的人来实现。我们使用下面的Operation来描述这个过程：
\begin{align}
  \rho \to \sum_i (M_i \otimes U_i) \rho_{AB} (M_i^\dagger \otimes
  U_i^\dagger) \quad \text{with} \quad \sum_i M_i^\dagger M_i = I_A
\end{align}

\bigskip
\hlr{LOCC作为Free Operations}

上面的第一种和第三种操作我们一般称之为LOCC（Local Operation and Classical
Communication）。我们发现LOCC操作不会产生entanglement，因此LOCC操作是entanglement
Resource Theory的Free Operations。

\bigskip
\hlr{Free States for Entanglement}

我们可以发现，Free States应该是那些不能通过LOCC操作产生的状态。我们发现这些状态正好是separable states。
\thm{
  Separable States are Free States.

  一个state $ \rho_{AB} $是free state当且仅当它可以写作：
  \begin{align}
    \rho_{AB} = \sum_i p_i \rho_A^i \otimes \rho_B^i \quad
    \text{with} \quad p_i \geq 0, \quad \sum_i p_i = 1.
  \end{align}
}

\subsubsection{Bell State LOCC interconversion}

我们希望研究LOCC能够怎么样子把哪些态变成另外一些态。

\bigskip
\hlr{Conversion of Bell State}

我们考虑Bell State $ \ket{\phi^+} = \frac{1}{\sqrt{2}} (\ket{00} +
\ket{11}) $以及另一个态$
|\theta\rangle=\cos(\theta)|00\rangle+\sin(\theta)|11\rangle $。
为此我们可以通过一个POVM测量实现，假设我们存在一个POVM 表示为$ M_A = \{A_1^\dagger A_1,
A_2^\dagger A_2 \} $，对于这两个矩阵定义为：
\begin{align}
  A_1=
  \begin{pmatrix}\cos(\theta)&\\&\sin(\theta)
  \end{pmatrix},\quad A_2=
  \begin{pmatrix}\sin(\theta)&\\&\cos(\theta)
  \end{pmatrix}.
\end{align}
我们可以定义一个LOCC操作：
\begin{itemize}
  \item STEP 1: 对于两个qubit的第一个qubit(A系统)进行POVM测量$ M_A $。根据POVM测量后坍缩的$
    \rho_k = A_k \rho A_k^\dagger $我们知道:
    \begin{itemize}
      \item 测量到1结果，系统变为：$ (A_1\otimes
        I)|\phi^+\rangle=\frac{1}{\sqrt{2}}(\cos(\theta)|00\rangle+\sin(\theta)|11\rangle)=\frac{1}{\sqrt{2}}|\theta\rangle
        $
      \item 测量到2结果，系统变为：$ (A_2\otimes
        I)|\phi^+\rangle=\frac{1}{\sqrt{2}}(\sin(\theta)|00\rangle+\cos(\theta)|11\rangle)
        = \frac{1}{\sqrt{2}} X \otimes X \ket{\theta} $
    \end{itemize}
  \item STEP 2: 如果得到1结果，A告诉B不进行任何操作；如果得到2结果，A告诉B对B系统进行X操作。
\end{itemize}

\bigskip
\hlr{Any State from Bell State}

实际上我们可以通过LOCC把一个bell state $ \ket{\phi^+} $变成任意的两个qubit的pure state $
|\psi\rangle_{AB} $。
但是原则上我们不能通过LOCC把一个一般的entangled state变成bell state。
变换操作需要使用Schimdt Decomposition来进行分析。我们定义：
\defi{Schmidt Decomposition

  对于一个定义在$ AB $系统上的pure state $ |\psi\rangle_{AB} $，我们总可以找到两个正交归一基$
  \{|u_i\rangle_A\}, \{|v_i\rangle_B\} $使得：
  \begin{align}
    |\psi\rangle_{AB} = \sum_{i=1}^r \sqrt{\lambda_i} |u_i\rangle_A
    \otimes |v_i\rangle_B, \quad 1 \leq r \leq
    \min(\dim\mathcal{H}_A, \dim\mathcal{H}_B)
  \end{align}
  其中$ \lambda_i > 0 $并且$ \sum_i \lambda_i = 1 $。我们称$ \{\lambda_i\}
  $为Schmidt Coefficients，$ r $为Schmidt Rank。
}
对于一个2 qubit system来说我们总能够把任意的pure state写作：
\begin{align}
  |\psi\rangle_{AB} = \sqrt{\lambda_1} |u_1\rangle_A \otimes
  |v_1\rangle_B + \sqrt{\lambda_2} |u_2\rangle_A \otimes |v_2\rangle_B,
\end{align}
于是我们可以构建一个Unitary满足：
\begin{align}
  U_A |0\rangle_A = |u_1\rangle_A, \quad U_A |1\rangle_A = |u_2\rangle_A, \\
  U_B |0\rangle_B = |v_1\rangle_B, \quad U_B |1\rangle_B = |v_2\rangle_B.
\end{align}
于是$  U = U_A \otimes U_B  $是一个local unitary操作，从$ \ket{\theta} $变成$
|\psi\rangle_{AB} $。因此我们发现只要能够把$ \ket{\psi^+} $变成$ \ket{\theta}
$，就可以通过local unitary把$ \ket{\theta} $变成任意的2 qubit pure state。

\subsubsection{Pure State Conversion and Majorization Theorem}

我们希望最一般的研究两个pure state $ |\psi\rangle_{AB} $和$ |\phi\rangle_{AB}
$之间的LOCC转换关系。需要满足什么关系我们才能把一个态LOCC变成另一个态。为此首先我们需要讨论一个数学定理：

\bigskip
\hlr{Majorization Theory}

\defi{Majorization of vectors

  对于一个向量$ \vec{x} $我们使用$ x^{\downarrow} $表示把$ \vec{x}
  $的分量从大到小排序后的向量。对于两个向量$ \vec{x}, \vec{y} \in \mathbb{R}^n $，我们定义$
  \vec{x} $被$ \vec{y} $ majorize，记作$ \vec{x} \prec \vec{y} $，当且仅当:
  \begin{itemize}
    \item $ \sum_ix_i=\sum_iy_i. $
    \item 对于任意$ k=1,2,\ldots,n-1 $，都有$ \sum_{i=1}^k x_i^{\downarrow}
      \leq \sum_{i=1}^k y_i^{\downarrow}. $
  \end{itemize}
}
对于两个向量如果其数值之和一样，我们可以定义majorization来描述一个partial ordering
relation。根据这个定义我们会发现，所有向量其实都是处于：
\begin{align}
  (\frac{1}{N},\frac{1}{N},\cdots,\frac{1}{N})\prec\vec{p}\prec
  (1,0,0,\cdots,0).
\end{align}

\bigskip
\hlr{Nielsen's Majorization Theorem}

我们发现对于pure state来说，LOCC的转换关系正好可以通过majorization来描述。
\thm{Nielsen's Majorization Theorem

  对于两个pure state定义在$ AB $系统上$
  \left|\varphi\right\rangle_{AB},\left|\psi\right\rangle_{AB} $我们考虑$
  \rho_A = \operatorname{Tr}_B(\ketbra{\psi}{\psi}) $的本征值计作$ \lambda(\psi) $那么：
  \begin{align}
    \left(\left|\varphi\right\rangle_{AB}\overset{\mathrm{LOCC}}{\operatorname*{\operatorname*{\longmapsto}}}\left|\psi\right\rangle_{AB}\right)\Longleftrightarrow\lambda(\varphi)\prec\lambda(\psi).
  \end{align}
}
我们注意，对于一个pure state我们永远可以做Schmidt Decomposition:
\begin{align}
  |\psi\rangle_{AB}=\sum_{i=1}^r\sqrt{\lambda_i}\left|u_i\right\rangle_A\otimes\left|v_i\right\rangle_B
  , \quad 1\leq r\leq\min(\dim\mathcal{H}_A,\dim\mathcal{H}_B)
\end{align}
因此对于A系统还是B系统做partial trace我们都会得到相同的本征值分布。因此这个定理告诉我们：
\begin{itemize}
  \item Tr掉一个系统的Majorization决定了LOCC的可转换性
\end{itemize}
我们之前讨论的$ \ket{\phi^+} $如果Tr掉A系统得到的本征值分布是$ (1/2,1/2) $，而任意两个qubit的pure
state Tr掉A系统得到的本征值分布都是$ (p,1-p) $，因此我们发现$ (1/2,1/2)\prec(p,1-p)
$对于任意$ p\in[0,1] $都成立，因此LOCC可以把$ \ket{\phi^+} $变成任意两个qubit的pure state。

\subsubsection{Pure State Asymptotic Conversion}

\bigskip
\hlr{Interconversion Rate}

\thm{
  Asymptotic Interconversion Rate

  对于任意的一个2 qubit pure state $ \ket{\varphi}_{AB} $我们发现下面两个态可以通过LOCC互相变换：
  \begin{align}
    |\varphi\rangle\langle\varphi|^{\otimes
    N}\overset{\mathrm{LOCC}}{\longleftrightarrow}\ketbra{\phi^+}{\phi^+}^{\otimes\left(S\left(\rho_{\varphi_{A}}\right)N\right)},
  \end{align}
  其中$ \rho_{\varphi_A} =
  \operatorname{Tr}_B(\ketbra{\varphi}{\varphi}) $，$ S(\rho) =
  -\operatorname{Tr}(\rho\log_2\rho) $是von Neumann entropy。
}
这告诉我们，我们可以通过LOCC把大量的普通纠缠态变成更少一点的bell state。反之，bell state也可以通过LOCC变成更大量的普通纠缠态。

\bigskip
\hlr{Interconversion Rate 的证明}

\begin{itemize}
  \item STEP 1: Unitary Transformation to Schmidt Form
\end{itemize}
对于任何的态我们都可以进行一个Schmidt Decomposition然后进行一个Local Unitary把态变成：
\begin{align}
  |\varphi\rangle=\sqrt{1-\lambda}|00\rangle+\sqrt{\lambda}|11\rangle.
\end{align}
因此对于$ N $个这样的态我们可以通过LOCC得到：
\begin{align}
  |\varphi\rangle^{\otimes
  N}=\left(\sqrt{1-\lambda}|00\rangle_{AB}+\sqrt{\lambda}|11\rangle_{AB}\right)^{\otimes
  N}.
\end{align}
如果我们把每一项写开会得到：
\begin{align}\label{eq:lec11-1}
  \left|\varphi\right\rangle^{\otimes
  N}=\sum_{k}\left(\sqrt{1-\lambda}\right)^{N-k}\left(\sqrt{\lambda}\right)^{k}\sum_{\pi}\left|\pi\left(\underbrace{1\cdots1}_{k}\underbrace{0\cdots0}_{N-k}\right)\right\rangle_{A}\left|\pi\left(\underbrace{1\cdots1}_{k}\underbrace{0\cdots0}_{N-k}\right)\right\rangle_{B}.
\end{align}
注意！A和B的$ \pi $需要是完全一样的，因为这两个qubit是完全纠缠的。
\begin{itemize}
  \item STEP 2: Measure the number of 1s
\end{itemize}
下面我们对于A系统进行一个测量，测量的内容是A系统中1的个数。对于这样的系统我们知道测量到$ k $个1，对于$ k = \lambda
N $的概率是最大的！因此我们以最大概率会坍缩到下面的态：
\begin{align}
  \left|\varphi\right\rangle^{\otimes N}\to\frac{1}{\binom{N}{\lambda
  N}^{1/2}}\sum_\pi\left|\pi\left(\underbrace{1\cdots1}_{\lambda
  N}0\cdots0\right)\right\rangle_A\left|\pi\left(\underbrace{1\cdots1}_{\lambda
  N}0\cdots0\right)\right\rangle_B.
\end{align}
我们计算右边线性叠加的态有多少个，也就是组合数$ \binom{N}{\lambda N} $。根据Von Noemann
Entropy的定义我们知道 $ \lambda=S(\varphi_A) $，因此我们有：
\begin{align}
  \binom{N}{\lambda N}\approx2^{N S{\rho_{\varphi_A}}},
\end{align}
因此我们这样的测量保证了得到了$ 2^{N S(\rho_{\varphi_A})} $个量子态的线性叠加。
\begin{itemize}
  \item STEP 3: Special Unitary to Bell Pairs
\end{itemize}
下面我们对于AB系统同时进行一个unitary变换把上面的态变成$ S(\rho_{\varphi_A})N $个bell
pairs。我们知道N个bell State的张量积的结果是所有可能的$ 2^N $个qubit的线性叠加：
\begin{align}
  \left|\psi^{+}\right\rangle^{\otimes
  N}=\frac{1}{\sqrt{2^N}}\sum_{j=0}^{2^N-1}\left|j\right\rangle_A\left|j\right\rangle_B.
\end{align}
其中$ j $是一个N位的二进制数。
因此我们可以构建一种Unitary同时作用在A和B系统上保证一种编码方式使得每一种permutation编码为前$
S(\rho_{\varphi_A})N $个qubit的二进制数。而后面的$ (1-S(\rho_{\varphi_A}))N
$个qubit全部变成$ |0\rangle $。这样我们就得到了$ S(\rho_{\varphi_A})N $个bell pairs。具体的编码方式如下：
\begin{enumerate}
  \item 后$ (1-\lambda) N $个qubit都是$ |0\rangle $，因为编码不需要他们
  \item 前$ \lambda N $个qubit，对应的二进制数$ j $从$ 0 $到$
    2^{S(\rho_{\varphi_A})N}-1 $进行编码每个对应一种permutation的构型。
\end{enumerate}
于是我们可以把之前的态通过LOCC变成$ S(\rho_{\varphi_A})N $个bell pairs:
\begin{align}
  \ket{\varphi}^{\otimes N} \to
  &\frac{1}{2^{S(\rho_A)N}}\sum_{j=1}^{2^{S(\rho)N}}|jj\rangle_{AB}\otimes|00\rangle_{AB}^{\otimes(1-S(\rho_A))N}\\
  = &\left|\varphi^+\right\rangle_{AB}^{\otimes
  S(\rho_A)N}\otimes|00\rangle_{AB}^{\otimes(1-S(\rho_A))N}.
\end{align}
\begin{itemize}
  \item STEP 4: Trace out the extra qubits
\end{itemize}
我们知道最终结果是一个product state，因此我们可以直接trace掉多余的qubit，得到最终结果：
\begin{align}
  \ket{\varphi}^{\otimes N} \to
  \left|\varphi^+\right\rangle_{AB}^{\otimes S(\rho_A)N}.
\end{align}
